% p151_2; source PDF page 151; manually verified editable table.
\begin{tikzpicture}[x=1cm,y=-1cm,inner sep=0pt,outer sep=0pt]
\path[use as bounding box] (0,0) rectangle (16.000,11.275);
\draw[line width=.4pt] (0.000,0.000) rectangle (1.600,0.855);
\node[align=center,text width=1.340cm,font={\linespread{1}\fontsize{10}{13.50}\selectfont\bfseries}] at (0.800,0.427) {分类};
\draw[line width=.4pt] (1.600,0.000) rectangle (7.800,0.855);
\node[align=center,text width=5.940cm,font={\linespread{1}\fontsize{10}{13.50}\selectfont\bfseries}] at (4.700,0.427) {案件事实};
\draw[line width=.4pt] (7.800,0.000) rectangle (10.000,0.855);
\node[align=center,text width=1.940cm,font={\linespread{1}\fontsize{10}{13.50}\selectfont\bfseries}] at (8.900,0.427) {结论};
\draw[line width=.4pt] (10.000,0.000) rectangle (16.000,0.855);
\node[align=center,text width=5.740cm,font={\linespread{1}\fontsize{10}{13.50}\selectfont\bfseries}] at (13.000,0.427) {称谓};
\draw[line width=.4pt] (0.000,0.855) rectangle (1.600,7.765);
\node[align=center,text width=1.340cm,font={\linespread{1}\fontsize{10}{13.50}\selectfont}] at (0.800,4.310) {帮助};
\draw[line width=.4pt] (1.600,0.855) rectangle (7.800,2.560);
\node[align=left,text width=5.940cm,font={\linespread{1}\fontsize{10}{13.50}\selectfont}] at (4.700,1.708) {提供不可能有用的帮助行为};
\draw[line width=.4pt] (7.800,0.855) rectangle (10.000,2.560);
\node[align=center,text width=1.940cm,font={\linespread{1}\fontsize{10}{13.50}\selectfont}] at (8.900,1.708) {帮助者无罪};
\draw[line width=.4pt] (10.000,0.855) rectangle (16.000,2.560);
\node[align=left,text width=5.740cm,font={\linespread{1}\fontsize{10}{13.50}\selectfont}] at (13.000,1.708) {对该案件事实，可称为“帮助未遂”，但不能称为“帮助犯未遂”“未遂的帮助犯”};
\draw[line width=.4pt] (1.600,2.560) rectangle (7.800,7.765);
\node[align=left,text width=5.940cm,font={\linespread{1}\fontsize{10}{13.50}\selectfont}] at (4.700,5.163) {提供可能有用的帮助行为。包括两种情形：（1）促进正犯的实行行为，但没有促进正犯的既遂结果。（甲拿着乙提供有用的钥匙盗窃，操作不当打不开门，用其他办法盗窃既遂）（2）促进正犯的实行行为，但正犯没有既遂（甲用乙的钥匙打开了门，盗窃时被抓）};
\draw[line width=.4pt] (7.800,2.560) rectangle (10.000,7.765);
\node[align=left,text width=1.940cm,font={\linespread{1}\fontsize{10}{13.50}\selectfont}] at (8.900,5.163) {帮助者成立帮助犯，但构成未遂};
\draw[line width=.4pt] (10.000,2.560) rectangle (16.000,7.765);
\node[align=left,text width=5.740cm,font={\linespread{1}\fontsize{10}{13.50}\selectfont}] at (13.000,5.163) {第一，若侧重于表达帮助犯的犯罪形态，对该案件事实可称为“帮助犯未遂”。第二，若侧重于表达帮助犯帮正犯帮到哪个阶段和程度，对该案件事实可称为“未遂的帮助犯”，意指没有帮到正犯既遂的程度，只帮到正犯的实行阶段（也称为未遂阶段）};
\draw[line width=.4pt] (0.000,7.765) rectangle (1.600,11.275);
\node[align=center,text width=1.340cm,font={\linespread{1}\fontsize{10}{13.50}\selectfont}] at (0.800,9.520) {教唆};
\draw[line width=.4pt] (1.600,7.765) rectangle (7.800,9.470);
\node[align=left,text width=5.940cm,font={\linespread{1}\fontsize{10}{13.50}\selectfont}] at (4.700,8.617) {教唆实施不可能有危险的行为};
\draw[line width=.4pt] (7.800,7.765) rectangle (10.000,9.470);
\node[align=center,text width=1.940cm,font={\linespread{1}\fontsize{10}{13.50}\selectfont}] at (8.900,8.617) {教唆者无罪};
\draw[line width=.4pt] (10.000,7.765) rectangle (16.000,9.470);
\node[align=left,text width=5.740cm,font={\linespread{1}\fontsize{10}{13.50}\selectfont}] at (13.000,8.617) {对该案件事实，可称为“未遂的教唆”，但不能称为“未遂的教唆犯”};
\draw[line width=.4pt] (1.600,9.470) rectangle (7.800,11.275);
\node[align=left,text width=5.940cm,font={\linespread{1}\fontsize{10}{13.50}\selectfont}] at (4.700,10.373) {教唆实施可能有危险的行为。正犯构成未遂};
\draw[line width=.4pt] (7.800,9.470) rectangle (10.000,11.275);
\node[align=left,text width=1.940cm,font={\linespread{1}\fontsize{10}{13.50}\selectfont}] at (8.900,10.373) {教唆者成立教唆犯，但构成未遂};
\draw[line width=.4pt] (10.000,9.470) rectangle (16.000,11.275);
\node[align=left,text width=5.740cm,font={\linespread{1}\fontsize{10}{13.50}\selectfont}] at (13.000,10.373) {对该案件事实，可称为“教唆犯未遂”，或“未遂的教唆犯”};
\end{tikzpicture}
